% ============================================================
%  ch01c  §1.3 整体几何朗兰兹对应
%  原文对应：ICM-arXiv.tex 750--780 行
%  编号：猜想 1.3.1
% ============================================================

\subsection{整体几何朗兰兹对应}\label{Global GLC}

如本文开头所述，（整体）几何朗兰兹纲领源自 Drinfeld 对函数域上 $\GL_2$ 朗兰兹猜想的证明。这一学科的早期发展，大多集中于对整体函数域 $F$ 的 Galois 表示（等价地，光滑代数曲线 $X$ 上的局部系统）构造 Hecke 特征层，参见 \cite{Drinfeld-GL2, Laumon-Geometric-Langlands, FGV-Geometric-Langlands-GLn}。

此后，在 \cite{Beilinson-Drinfeld-Hitchin} 的工作之后，整个纲领的范围发生了转变：Beilinson--Drinfeld 在其中给出了整体几何朗兰兹对应的一个粗略范畴式表述。这一表述随后由 Arinkin--Gaitsgory 在 \cite{Arinkin-Gaitsgory-Singular-Support} 中精确化，下面加以回顾。

设 $X$ 是 $F=\bC$ 上的光滑射影曲线。在自守一侧，令 $\mathbf{D}_c(\Bun_G)$ 为 $X$ 上主 $G$-丛的模叠 $\Bun_G$ 上凝聚 D-模构成的 $\infty$-范畴。在 Galois 一侧，令 $\Coh(\Loc_{\hat{G}})$ 为 $X$ 上 de Rham $\hat{G}$-局部系统（又称带平坦联络的主 $\hat{G}$-丛）的模叠 $\Loc_{\hat{G}}$ 上凝聚层构成的 $\infty$-范畴。

\begin{conj}\label{Conj: GGLC}
存在 $\infty$-范畴之间的典范等价
\[
\bL_G:  \Coh(\Loc_{\hat{G}})\cong \mathbf{D}_c(\Bun_G),
\]
且满足一列自然的相容性条件。
\end{conj}

我们只简述其中最重要的相容性条件，其余见 \cite{Arinkin-Gaitsgory-Singular-Support}。注意两侧都带有一族相互交换的算子的作用，这些算子由 $x\in X$ 与 $V\in \Coh(\bB\hat{G}_\bC)^\heart$ 标号。具体地说，对每一点 $x\in X$ 有求值映射 $\Loc_{\hat{G}}\to \bB\hat{G}_\bC$，于是每个 $V\in \Coh(\bB\hat{G}_\bC)^\heart$ 通过拉回给出 $\Loc_{\hat{G}}$ 上的向量丛 $\widetilde{V}_x$，后者再以张量积作用于 $\Coh(\Loc_{\hat{G}})$。另一方面，有 Hecke 算子 $H_{V,x}$ 作用于 $\mathbf{D}_c(\Bun_G)$，其做法是借助（几何 Satake \eqref{factorizable Sat} 的 D-模版本）把层 $\Sat_{\{1\}}(V)|_x$ 作卷积。等价 $\bL_G$ 应当把这些算子的作用相互对应起来。

尽管这一猜想仍远未解决，人们已知 $\Loc_{\hat{G}}$ 上完全复形的范畴 $\Perf(\Loc_{\hat{G}})$ 作用于 $\mathbf{D}_c(\Bun_G)$，该作用通常称为谱作用；在此作用下，$\widetilde{V}_x\in \Perf(\Loc_{\hat{G}})$ 对 $\mathbf{D}_c(\Bun_G)$ 的作用由 Hecke 算子 $H_{V,x}$ 给出。

如今猜想 \ref{Conj: GGLC} 有时被称为整体几何朗兰兹猜想的 de Rham 版本，因为这类猜想式等价还有其他版本，下面略作介绍。

首先，本着非交换 Hodge 理论的精神，猜想 \ref{Conj: GGLC} 应当存在一个经典极限，有时称为整体几何朗兰兹的 Dolbeault 版本。其精确表述（据作者所知）尚不清楚，但一般说来，它归结为 $G$ 与 $\hat{G}$ 的 Hitchin 纤维化的对偶（在镜像对称的意义下）；这一点已在 \cite{Donagi-Pantev-Hitchin, Chen-Zhu-Geometric-Langlands-char-p} 中“一般地”建立。通过在正特征下扭变／形变这种对偶，可以证明猜想 \ref{Conj: GGLC} 的特征 $p$ 类比（当然也只是“一般地”，见 \cite{Bezrukavnikov-Braverman-Geometric-Langlands-char-p, Chen-Zhu-Nonabelian-Hodge, Chen-Zhu-Geometric-Langlands-char-p}）。有意思的是，这个“一般”的特征 $p$ 等价可用来给出 \cite{Beilinson-Drinfeld-Hitchin} 主要结果的新证明（至少对 $G=\GL_n$，见 \cite{Bezrukavnikov-Travkin-Quantization}）。

\cite{Beilinson-Drinfeld-Hitchin} 的工作（从而整体几何朗兰兹的 de Rham 版本）深受共形场论的影响。另一方面，受拓扑场论启发，Ben-Zvi 与 Nadler \cite{Ben-Zvi-Nadler-Betti-Geometric-Langlands} 提出了猜想 \ref{Conj: GGLC} 的 Betti 版本：在自守一侧，把 $\Bun_G$ 上 $D$-模范畴替换为（$\Bun_G$ 解析化后的）$\bC$-向量空间层范畴；在 Galois 一侧，把 $\Loc_{\hat{G}}$ 替换为 Betti $\hat{G}$-局部系统的模空间（又称 $X$ 的基本群的 $\hat{G}$-值表示）。

Riemann--Hilbert 对应允许在 de Rham $\hat{G}$-局部系统与 Betti $\hat{G}$-局部系统之间过渡，但这一过渡是超越的。所以猜想 \ref{Conj: GGLC} 与其 Betti 类比并无直接联系。

最近，Arinkin 等人 \cite{AGKRRV} 提出了猜想 \ref{Conj: GGLC} 的又一个版本，它直接把 de Rham 版本与 Betti 版本联系起来，同时又包含用 $\ell$-进层表述的版本。因此它与函数域上的经典朗兰兹对应关系更密切，这将在 \S\ref{S: Global Langlands function fields} 中讨论。
